\chapter{Joins relacionales y su conexión con grafos}
\label{cap:joins-y-grafos}

\begin{intuicion}
Hay dos formas equivalentes de mirar el mismo problema:
\begin{itemize}
    \item ``Subgraph matching sobre un grafo etiquetado.''
    \item ``Multi-way join sobre tablas relacionales.''
\end{itemize}
La equivalencia importa porque la comunidad de bases de datos
estuvo trabajando \emph{décadas} en la teoría de joins, y en ese
camino derivó el \emph{AGM bound} ---una cota inferior sobre cuán
chico puede ser cualquier algoritmo correcto--- y los \emph{worst-case
optimal joins} (WCOJ), algoritmos que alcanzan el bound. Este
capítulo arranca explicando \emph{qué} es un join (sin asumir que
sabés SQL), construye la equivalencia ``subgraph matching = join'',
deriva el AGM bound desde cero con dos ejemplos chicos, y motiva
\emph{por qué} el motor incluye un operador \emph{worst-case-optimal}
para \(K_4\).
\end{intuicion}

\section{Joins desde cero}

\subsection{Relaciones}

\begin{definition}[Relación]
Una \emph{relación} \(R\) sobre los atributos
\(A_1, A_2, \dots, A_d\) es un conjunto finito de tuplas, donde cada
tupla \(t \in R\) asigna un valor a cada atributo. Notación:
\(R(A_1, \dots, A_d)\). El número \(d\) se llama \emph{arity} de \(R\).
\end{definition}

Una relación es básicamente una \emph{tabla}: las columnas son los
atributos, las filas son las tuplas. Sin claves duplicadas, sin orden
entre filas. Por ejemplo:

\[
\begin{array}{l}
\texttt{Auth}(\texttt{user}, \texttt{host}) = \\[0.2em]
\quad\{ (\texttt{alice}, \texttt{db-01}),\ (\texttt{bob}, \texttt{web-02}),\ (\texttt{alice}, \texttt{web-02}) \}
\end{array}
\]

\subsection{Join binario natural}

\begin{definition}[Join natural binario]
Dadas dos relaciones \(R(A_1, \dots, A_d)\) y \(S(B_1, \dots, B_e)\),
sea \(C = \{A_1, \dots, A_d\} \cap \{B_1, \dots, B_e\}\) el conjunto
de atributos compartidos. El \emph{join natural} \(R \bowtie S\) es
la relación sobre \(\{A_1, \dots, A_d\} \cup \{B_1, \dots, B_e\}\)
cuyas tuplas son:
\[
R \bowtie S = \{ t \cup s : t \in R, s \in S, \forall c \in C.\ t[c] = s[c] \}
\]
donde \(t \cup s\) es la tupla que combina los valores de \(t\) y
\(s\) (consistentes en los atributos compartidos).
\end{definition}

\begin{ejemplo}
Sean
\[
\begin{array}{l}
R(\texttt{user}, \texttt{host}) = \{ (\texttt{alice}, \texttt{db-01}),\ (\texttt{bob}, \texttt{web-02}) \} \\[0.3em]
S(\texttt{host}, \texttt{ip}) = \{ (\texttt{db-01}, \texttt{10.0.0.5}),\ (\texttt{web-02}, \texttt{8.8.8.8}),\ (\texttt{web-02}, \texttt{1.1.1.1}) \}
\end{array}
\]

El atributo compartido es \texttt{host}. Entonces:

\[
\begin{array}{l}
R \bowtie S = \{ \\
\quad (\texttt{alice}, \texttt{db-01}, \texttt{10.0.0.5}), \\
\quad (\texttt{bob}, \texttt{web-02}, \texttt{8.8.8.8}), \\
\quad (\texttt{bob}, \texttt{web-02}, \texttt{1.1.1.1}) \\
\}
\end{array}
\]

\noindent
La tupla \((\texttt{alice}, \texttt{web-02})\) de \(R\) NO matchea
con nada de \(S\) ---porque \(R\) no la contiene---. La tupla
\((\texttt{bob}, \texttt{web-02})\) sí, y matchea con dos filas de
\(S\), generando dos tuplas en la salida.
\end{ejemplo}

\subsection{Multi-way join}

\begin{definition}[Multi-way join]
Dadas \(n\) relaciones \(R_1, R_2, \dots, R_n\), el \emph{multi-way
join} \(R_1 \bowtie R_2 \bowtie \dots \bowtie R_n\) es:
\[
\bigcup \{ t_1 \cup t_2 \cup \dots \cup t_n :
    t_i \in R_i,\ \text{los \(t_i\) son consistentes en cada atributo compartido} \}
\]
\end{definition}

\noindent
El join natural es asociativo y conmutativo, así que el orden de
ejecución no cambia el resultado, pero sí cambia el costo: ejecutar
\((R_1 \bowtie R_2) \bowtie R_3\) puede ser muy distinto de
\((R_1 \bowtie R_3) \bowtie R_2\) en términos de \emph{rows
intermedias}. Esa es la pieza que el \emph{query planner} de un SIEM
clásico optimiza, y donde los \emph{worst-case-optimal joins}
cambian las reglas.

\section{Equivalencia: subgraph matching es multi-way join}

Ahora la pieza fundamental del capítulo.

\begin{theorem}[Subgraph matching como join]
Sea \(P\) un patrón con vértices \((p_0, \dots, p_k)\) y aristas
\(\{(p_{i_1}, p_{j_1}), \dots, (p_{i_m}, p_{j_m})\}\). Sea \(G\) un
grafo. Para cada arista \((p_{i_l}, p_{j_l})\) del patrón con
etiqueta \(L_l\), definamos la relación:
\[
R_l(\texttt{src}_{i_l}, \texttt{dst}_{j_l}) =
\{(u, v) : (u, v) \in \E(G),\ \mathrm{type}_E((u, v)) = L_l \}
\]
Entonces el conjunto de embeddings homomórficos de \(P\) en \(G\) es
exactamente:
\[
R_1 \bowtie R_2 \bowtie \dots \bowtie R_m
\]
donde el atributo \(\texttt{src}_i\) (resp.\ \(\texttt{dst}_j\))
identifica el vértice \(p_i\) del patrón.
\end{theorem}

\begin{proof}[Idea de la prueba]
Una asignación \(\sigma\) que valga como embedding tiene que cumplir
dos cosas: (a) cada arista del patrón existe en \(G\) con la etiqueta
correcta ---eso es exactamente decir ``\((\sigma(p_i), \sigma(p_j))
\in R_l\)''---; (b) la asignación es consistente: si dos aristas del
patrón comparten un vértice (digamos \(p_i\) aparece en
\((p_i, p_j)\) y en \((p_i, p_k)\)), entonces el valor que \(\sigma\)
le asigna a \(p_i\) es el mismo en ambas aristas. La consistencia es
exactamente lo que el join natural impone vía atributos compartidos.
\(\square\)
\end{proof}

\begin{ejemplo}[\(K_4\) como join de 6 relaciones]
El patrón \(K_4\) tiene cuatro vértices \(p_0, p_1, p_2, p_3\) y las
seis aristas dirigidas entre cada par:
\[
\{(p_0, p_1), (p_0, p_2), (p_0, p_3),
  (p_1, p_2), (p_1, p_3), (p_2, p_3)\}
\]
El join correspondiente es:
\[
R_{01}(\texttt{src}_0, \texttt{dst}_1)
  \bowtie R_{02}(\texttt{src}_0, \texttt{dst}_2)
  \bowtie R_{03}(\texttt{src}_0, \texttt{dst}_3)
  \bowtie R_{12}(\texttt{src}_1, \texttt{dst}_2)
  \bowtie R_{13}(\texttt{src}_1, \texttt{dst}_3)
  \bowtie R_{23}(\texttt{src}_2, \texttt{dst}_3)
\]

Si todas las aristas del grafo tienen la misma etiqueta, las seis
relaciones son \emph{la misma}: \(R_l = \E(G)\) para todo \(l\). El
join se vuelve un \emph{self-join} de \(\E\) consigo misma seis
veces, con los atributos renombrados.
\end{ejemplo}

\subsection{Por qué la equivalencia importa}

Dos consecuencias prácticas:

\begin{enumerate}
    \item Cualquier teorema sobre joins relacionales se aplica a
    subgraph matching. En particular: el \emph{AGM bound}
    (próxima sección) y los \emph{worst-case-optimal joins} (LFTJ).
    \item Cualquier algoritmo de join se puede usar como matcher.
    En particular: \emph{semi-join reducer} de Yannakakis para
    consultas \(\alpha\)-acíclicas (capítulo~\ref{cap:yannakakis-chain}).
\end{enumerate}

\section{El AGM bound}

\subsection{Motivación}

Antes de enunciarlo, una pregunta concreta: para un join de tres
relaciones, cada una con \(m\) tuplas, ¿cuántas tuplas \emph{en el
peor caso} puede tener la salida?

\begin{itemize}
    \item Si las tres son joinadas en cadena (atributos compartidos
    formando una secuencia \(A \to B \to C \to D\)), la salida puede
    crecer hasta \(m^3\): cada tupla del primer join se replica para
    cada coincidencia en el segundo.
    \item Si las tres forman un \emph{triángulo} (compartiendo
    atributos en forma cíclica), \emph{lo más} que la salida puede
    crecer es \(m^{1.5}\). \emph{No} \(m^3\). Esa es la sorpresa.
\end{itemize}

¿Por qué? Porque un triángulo tiene una \emph{estructura} más rígida
que una cadena: las tres relaciones se restringen mutuamente. Esta
intuición se hace precisa con el \emph{fractional edge cover number}.

\subsection{Cobertura fraccionaria}

\begin{definition}[Hipergrafo de la consulta]
Dado un join \(R_1 \bowtie \dots \bowtie R_m\) sobre atributos
\(\mathcal{A}\), su \emph{hipergrafo} \(H(Q) = (\mathcal{A}, F)\)
tiene como vértices los atributos y como hiperaristas las relaciones:
cada \(R_l\) corresponde a la hiperarista
\(F_l = \{A : A\text{ es un atributo de } R_l\}\).
\end{definition}

\begin{definition}[Edge cover fraccionario]
Una \emph{cobertura fraccionaria} de \(H(Q)\) es una asignación de
pesos \(x_l \geq 0\), uno por hiperarista \(F_l\), tal que para todo
atributo \(A \in \mathcal{A}\):
\[
\sum_{l : A \in F_l} x_l \geq 1.
\]
El \emph{fractional edge cover number} de \(Q\) es:
\[
\rhostar(Q) = \min \sum_{l=1}^{m} x_l, \quad \text{sujeto a las restricciones de cobertura.}
\]
\end{definition}

En palabras: cada atributo tiene que estar ``cubierto'' por las
hiperaristas que lo contienen, con peso total \(\geq 1\). El
\(\rhostar\) es la mínima suma de pesos necesaria.

\begin{ejemplo}[Triángulo]
Para el join triangular
\(R(A, B) \bowtie S(B, C) \bowtie T(A, C)\):
\begin{itemize}
    \item Hipergrafo: vértices \(\{A, B, C\}\); hiperaristas
    \(F_R = \{A, B\}\), \(F_S = \{B, C\}\), \(F_T = \{A, C\}\).
    \item Cobertura fraccionaria: asignar \(x_R = x_S = x_T = 0.5\)
    cubre cada atributo con peso \(0.5 + 0.5 = 1\). La suma total es
    \(1.5\).
    \item Cualquier cobertura con menos suma viola alguna restricción.
    Por lo tanto \(\rhostar = 1.5\).
\end{itemize}
\end{ejemplo}

\begin{ejemplo}[Cadena]
Para el join en cadena
\(R(A, B) \bowtie S(B, C) \bowtie T(C, D)\):
\begin{itemize}
    \item Atributo \(A\): sólo aparece en \(R\), entonces
    \(x_R \geq 1\).
    \item Atributo \(D\): sólo aparece en \(T\), entonces
    \(x_T \geq 1\).
    \item Con \(x_R = x_T = 1\), los atributos \(B\) (en \(R\) y
    \(S\)) y \(C\) (en \(S\) y \(T\)) ya están cubiertos por
    \(R\) y \(T\) respectivamente; \(x_S = 0\) basta.
    \item Suma: \(1 + 0 + 1 = 2\), pero podríamos pensar que \(S\)
    necesita estar cubierta. No: la \emph{cobertura} es de
    \emph{atributos}, no de aristas. \(\rhostar = 2\).
\end{itemize}

Wait. Eso no se condice con lo que dijimos antes (``cadena de tres
relaciones, salida puede ser \(m^3\)''). Veamos: cadena de tres
relaciones sobre cuatro atributos, sí, \(\rhostar = 2\). Pero la
salida es \emph{a lo sumo} \(m^2\) bajo el AGM bound. La intuición
de ``\(m^3\)'' que tenía uno antes era una sobre-estimación basada
en el caso ``join naive''; el AGM bound nos dice que en realidad la
estructura del triángulo y la cadena no son tan distintas como
pensábamos. La verdadera diferencia entre cadena y triángulo es de
\(m^2\) versus \(m^{1.5}\); ambas son menores que \(m^3\).
\end{ejemplo}

\subsection{El teorema}

\begin{theorem}[Atserias-Grohe-Marx 2008]
Para todo join \(Q = R_1 \bowtie \dots \bowtie R_m\) y toda
asignación de tamaños \(|R_l| \leq m\):
\[
|Q| \leq m^{\rhostar(Q)}.
\]
Además, esta cota es \emph{ajustada}: existen instancias donde
\(|Q|\) alcanza \(m^{\rhostar(Q)}\) hasta factores constantes.
\end{theorem}

\begin{proof}[Idea de la prueba]
La parte ``superior'' (\(|Q| \leq m^{\rhostar}\)) sale de aplicar la
desigualdad de Shearer (información mutua sub-aditiva sobre
particiones del producto) a la entropía de la distribución uniforme
sobre \(Q\). La parte ``ajustada'' se demuestra construyendo
instancias específicas (típicamente, todos los \(R_l\) son
\(m^{1/\rhostar}\) copias del Cartesiano más un balanceo). \(\square\)
\end{proof}

\begin{nota}
Lo crucial: \(\rhostar\) depende \emph{sólo de la estructura del
patrón}, no del tamaño de los datos. Es un invariante del problema
que dice ``este join puede crecer a lo sumo así''.
\end{nota}

\subsection{Worst-case optimal joins}

\begin{definition}[Algoritmo WCOJ]
Un algoritmo de join es \emph{worst-case-optimal} sii su tiempo de
ejecución es \(\widetilde{O}(m^{\rhostar(Q)})\) (el \(\widetilde{O}\)
oculta factores logarítmicos).
\end{definition}

\noindent
Los algoritmos clásicos de join (\emph{nested-loop},
\emph{hash-join}, \emph{sort-merge}) ejecutados \emph{de a dos
relaciones a la vez} \textbf{NO son WCOJ} para hipergrafos cíclicos
como el triángulo: pueden producir resultados intermedios de tamaño
\(m^2\) aun cuando la salida final es \(m^{1.5}\).

\begin{itemize}
    \item Triángulo: \(\rhostar = 1.5\). Hash-join binario es
    \(O(m \cdot m) = O(m^2)\) en el peor caso. WCOJ es \(O(m^{1.5})\).
    Ratio: \(\sqrt{m}\). Para \(m = 10^6\), eso es 1000\(\times\)
    más rápido.
    \item \(K_4\): \(\rhostar = 2\). Hash-join secuencial es
    \(O(m^3)\) en el peor caso. WCOJ es \(O(m^2)\). Para \(m = 10^4\),
    eso es \(10^4\times\) más rápido.
\end{itemize}

\noindent
El primer algoritmo WCOJ fue NPRR (Ngo-Porat-Re-Rudra, 2012). LFTJ
(Veldhuizen, 2014) ---el que GraphHunter implementa para \(K_4\)---
es una variante más implementable. El capítulo~\ref{cap:lftj} lo
desarrolla.

\section{Por qué entonces no toda hipótesis usa WCOJ}

Pregunta natural: si WCOJ es estrictamente mejor en el peor caso,
¿por qué el motor cae al DFS por defecto y sólo usa LFTJ para \(K_4\)
y Yannakakis para \(\alpha\)-acíclicos?

Tres razones:

\begin{enumerate}
    \item \textbf{El bound es de peor caso, no de caso típico.} En
    grafos ``reales'' (\emph{scale-free}, etiquetas que filtran
    masivamente, ventana temporal angosta), el DFS bien-podado puede
    ser ya \(O(m \cdot k)\) o mejor ---\emph{ganando} al WCOJ por
    constantes y por menor sobrecarga.
    \item \textbf{Implementación.} LFTJ requiere un \emph{trie}
    construido sobre las aristas, indexado consistentemente, con
    \emph{multi-way intersection} SIMD. Construirlo es trabajo;
    mantenerlo bajo ingestas streaming es más trabajo. Para
    hipótesis donde el DFS ya alcanza, no vale la pena.
    \item \textbf{Generalidad.} El DFS funciona para
    \emph{cualquier} hipótesis (chain, graph, con predicados
    arbitrarios). LFTJ está especializado a la \emph{shape}
    estructural; agregar predicados o etiquetas heterogéneas requiere
    extensiones.
\end{enumerate}

El \emph{planner} (capítulo~\ref{cap:planner}) materializa este
trade-off: dispatcha al WCOJ cuando \emph{sabe} que la estructura lo
permite, cae al DFS en caso contrario.

\section{Resumen}

\begin{itemize}
    \item Un \textbf{join natural} es la operación que junta dos
    tablas matcheando los valores en sus columnas compartidas.
    Multi-way join generaliza a \(n\) tablas.
    \item \textbf{Subgraph matching homomórfico es exactamente un
    multi-way join} de las relaciones de aristas, con consistencia
    impuesta vía atributos compartidos.
    \item El \textbf{AGM bound}
    (\(|Q| \leq m^{\rhostar(Q)}\)) acota el tamaño de la salida en
    función del \emph{fractional edge cover number} del hipergrafo
    de la consulta.
    \item \(\rhostar\) es \(1.5\) para triángulo, \(2\) para
    \(K_4\), \(2\) para cadena de 3 aristas, \(k\) para cadena de
    \(k\) aristas.
    \item Los \textbf{worst-case-optimal joins} (NPRR, LFTJ) alcanzan
    el AGM bound; los joins binarios clásicos pueden no alcanzarlo
    (sobre todo en hipergrafos cíclicos).
    \item GraphHunter usa LFTJ para \(K_4\) (capítulo~\ref{cap:lftj})
    y Yannakakis para \(\alpha\)-acíclicos (capítulos
    \ref{cap:yannakakis-chain} y \ref{cap:yannakakis-tree});
    para todo lo demás, DFS con \emph{prunes}.
\end{itemize}

\begin{ejercicios}
\begin{enumerate}
    \item \starred{} Calculá \(\rhostar\) para una \emph{estrella}
    de \(k\) brazos (un vértice central con \(k\) aristas salientes a
    vértices distintos). ¿Qué AGM bound da? ¿Es ajustado en grafos
    reales?

    \item Calculá \(\rhostar\) para el grafo bipartito completo
    \(K_{3,3}\) (visto como hipergrafo con 9 hiperaristas binarias).
    Comparalo con el bound trivial \(m^9\). ¿Qué dice esto sobre la
    posible eficiencia de un WCOJ para \(K_{3,3}\)?

    \item \starred{} Mostrá que para una cadena de \(k\) aristas
    el AGM bound es ajustado: construí un grafo de \(m\) aristas
    donde el join correspondiente tenga exactamente \(\Theta(m^k)\)
    matches.

    \item Diseñá un \emph{nested-loop join} naive para
    \(R \bowtie S\) (dos relaciones binarias) y demostrá que su
    complejidad es \(O(|R| \cdot |S|)\). ¿Bajo qué condiciones
    (sobre el atributo compartido) un \emph{hash-join} mejora a
    \(O(|R| + |S|)\)?

    \item Imaginá una hipótesis que el motor recibe con
    \emph{shape} \emph{chain} de longitud 4. Identificá la
    estructura como un join de 4 relaciones binarias y dibujá su
    hipergrafo. ¿Cuál es \(\rhostar\)?

    \item ¿Por qué el AGM bound no contradice el hecho de que
    subgraph isomorphism sea NP-completo? (Pista: ¿qué crece con
    qué?)
\end{enumerate}
\end{ejercicios}
